\chapter{Đừng Vay Để Nuôi Thói Quen}
\markboth{Đừng Vay Để Nuôi Thói Quen}{Do Not Borrow to Feed a Habit}

\section{Khoản Vay Tạm Và Cái Giá Dài Hơn Một Tháng}
\begin{truyen}{Khi Thiếu Tiền Là Nghĩ Ngay Đến Mượn}{Thinking of Borrowing Before Thinking of Cutting Back}
\chuhoa{C}{ó giai đoạn cứ gần cuối tháng là người thầy lại nghĩ tới chuyện xoay tạm một khoản. Ban đầu rất nhỏ, chỉ để lấp một chỗ hụt. Nhưng rồi việc ấy lặp lại thành một lối mòn trong đầu.}
\textit{(There was a period when, near the end of each month, the teacher would think about borrowing a small amount. At first it was tiny, just to fill a gap. But the act repeated until it became a groove in the mind.)}

Mỗi lần thiếu, thay vì nhìn lại những khoản chưa cần, anh lại tìm một người quen để mượn qua. Cái dễ của cách ấy làm người ta tưởng mình đã giải quyết được vấn đề. Nhưng thật ra, mình chỉ đẩy sự mệt mỏi sang tháng sau.
\textit{(Each time money ran short, instead of reexamining unnecessary expenses, he looked for someone to borrow from. The convenience of that method makes people think the problem is solved. In truth, the burden is merely pushed into the next month.)}

Đến lúc tiền chưa về mà khoản cần trả đã tới, anh mới thấy rõ: nợ không chỉ lấy mất tiền của tương lai, nó còn lấy cả sự thảnh thơi của hiện tại. Một người sống dựa vào vay tạm lâu ngày thường đánh mất khả năng tự sửa thói quen của mình.
\textit{(When repayment arrived before new income did, he finally saw clearly: debt does not only take future money, it also steals present peace. A person who relies too long on temporary borrowing often loses the ability to correct his own habits.)}
\end{truyen}

\begin{baihoc}
Khoản vay dễ nhất thường cũng là thứ làm ta chậm sửa mình nhất.
\textit{(The easiest loan is often what delays self-correction the longest.)}
\end{baihoc}

\begin{ghinhoanh}
Easy borrowing hides hard consequences.

Debt often postpones discipline.
\end{ghinhoanh}

\begin{tuhocnhanh}
1. Khi thiếu tiền, phản ứng đầu tiên của em là cắt bớt hay đi vay? \textit{When money is short, is your first response to cut back or to borrow?}

2. Em đang đẩy gánh nặng hôm nay sang tháng sau ở điểm nào? \textit{Where are you pushing today’s burden into the next month?}
\end{tuhocnhanh}
\ngancach

\section{Liệu Cơm Gắp Mắm}
\begin{truyen}{Sống Theo Sức Mình Là Trí Khôn Cũ Nhưng Không Bao Giờ Cũ}{Living Within Means Is Old Wisdom That Never Ages}
\chuhoa{N}{gười xưa nói ``liệu cơm gắp mắm''. Câu nói nghe mộc mạc, nhưng nó chứa cả một nền khôn ngoan về cách sống.}
\textit{(The old saying goes, ``Measure the fish sauce according to the rice.'' It sounds simple, yet it contains an entire wisdom about living.)}

Biết liệu sức mình không làm con người nhỏ đi. Trái lại, nó giúp người ta tránh được rất nhiều cảnh khó xử. Không phải thứ gì nhà người khác có thì nhà mình cũng phải có ngay. Không phải hoàn cảnh nào cũng đáng vay tiền để giữ mặt mũi.
\textit{(Knowing one’s means does not make a person smaller. On the contrary, it helps avoid many difficult situations. Not everything another household owns must be owned immediately in one’s own. Not every social situation is worth borrowing money just to save face.)}

Người thầy dần hiểu rằng sống vừa sức không phải dấu hiệu của thất bại. Nó là sự trưởng thành của người biết nhìn xa hơn một cuộc vui, một bữa gặp gỡ, một món đồ đang làm mình nao lòng.
\textit{(The teacher gradually understood that living within one’s means is not a sign of failure. It is the maturity of someone who sees farther than one gathering, one pleasure, or one object that currently stirs desire.)}
\end{truyen}

\begin{baihoc}
Sống vừa sức là một dạng tự trọng rất thực tế.
\textit{(Living within one’s means is a very practical form of self-respect.)}
\end{baihoc}

\begin{ghinhoanh}
Live within your means.

That is practical dignity.
\end{ghinhoanh}

\begin{tuhocnhanh}
1. Điều gì làm em khó chấp nhận giới hạn của mình nhất? \textit{What makes it hardest for you to accept your limits?}

2. Em có đang vay để giữ thể diện không? \textit{Are you borrowing to preserve appearances?}
\end{tuhocnhanh}
\ngancach

\section{Sửa Thói Quen Trước Khi Tìm Tiền}
\begin{truyen}{Cái Thiếu Không Chỉ Là Tiền}{The Lack Is Not Only Financial}
\chuhoa{S}{au nhiều lần loay hoay, người thầy bắt đầu nhận ra cái thiếu thật sự của mình không còn nằm hoàn toàn ở số tiền. Cái thiếu lớn hơn là khả năng nói ``không'' với những thứ chưa cần.}
\textit{(After many struggles, the teacher began to see that his true lack was no longer located entirely in the amount of money. The deeper lack was the ability to say ``no'' to what was not yet necessary.)}

Khi chưa sửa được chỗ đó, vay thêm chỉ giống như đổ nước vào một chiếc thùng thủng. Thùng đầy lên một chút rồi lại vơi. Người ta tưởng mình thiếu nước, nhưng thật ra cái đáng sửa đầu tiên là lỗ rò.
\textit{(Until that weakness was corrected, borrowing more was like pouring water into a leaking bucket. It filled slightly, then emptied again. One thinks the problem is not enough water, but the first thing to repair is actually the leak.)}

Từ lúc hiểu điều ấy, anh bớt hỏi ``vay ở đâu'' và bắt đầu hỏi ``mình đã phung phí ở đâu''. Câu hỏi sau khó nghe hơn, nhưng là cánh cửa duy nhất dẫn tới một đời sống bền hơn.
\textit{(Once he understood this, he asked less often, ``Where can I borrow?'' and more often, ``Where have I been wasteful?'' The second question is harder to hear, but it is the only door that leads to a more stable life.)}
\end{truyen}

\begin{baihoc}
Trước khi tìm thêm tiền, hãy tìm chỗ mình đang làm rò tiền.
\textit{(Before searching for more money, find where you are leaking the money you already have.)}
\end{baihoc}

\begin{ghinhoanh}
Fix the leak before adding water.

Self-honesty saves more than panic.
\end{ghinhoanh}

\begin{tuhocnhanh}
1. Lỗ rò tài chính lớn nhất của em là gì? \textit{What is your biggest financial leak?}

2. Em dám sửa nó từ ngày nào? \textit{From which day will you begin to repair it?}
\end{tuhocnhanh}

\begin{turanminh}
1. Tôi không được lấy nợ để che cho sự buông tay của mình trong chi tiêu.

2. Trước khi hỏi vay ở đâu, tôi phải tự hỏi mình đã phung phí ở đâu.

3. Nếu chưa sửa được thói quen tiêu tiền, mọi khoản vay chỉ làm cái mệt của tháng sau nặng hơn tháng này.
\end{turanminh}

\begin{dayhoc}
1. Với học sinh, bài học từ chương này là đừng tập thói quen sống trước khả năng của mình.

2. Có thể nói với các em rằng người không biết chịu giới hạn rất dễ mang nợ không chỉ bằng tiền, mà còn bằng áp lực và mất tự do.

3. Một câu hỏi phù hợp cho lớp là: ``Điều gì đáng sợ hơn, thiếu tiền tạm thời hay sống không biết điểm dừng?''
\end{dayhoc}